\chapter{Projekt i implementacja środowiska Ultimate Texas Hold’em}
\label{chap:srodowisko-uth}

W niniejszym rozdziale przedstawiono projekt oraz implementację środowiska
Ultimate Texas Hold’em wykorzystywanego w dalszej części pracy do uczenia
i oceny agentów. Opisane wcześniej zasady gry zostały odwzorowane w postaci
deterministycznego i testowalnego silnika, który zarządza przebiegiem rozdania,
udostępnia agentowi legalne akcje oraz rozlicza wynik finansowy rozgrywki.

Szczególną uwagę poświęcono zachowaniu niepełnej informacji charakterystycznej
dla pokera. Agent otrzymuje wyłącznie informacje, które byłyby dostępne
rzeczywistemu graczowi, natomiast karty krupiera, karty spalone oraz kolejność
pozostałych kart w talii pozostają elementami wewnętrznego stanu środowiska.
Takie rozdzielenie pozwala ograniczyć ryzyko niezamierzonego ujawnienia
informacji podczas uczenia modeli.

Implementację podzielono na niezależne moduły odpowiedzialne między innymi
za reprezentację kart, ocenę układów, sterowanie przebiegiem rozdania,
walidację akcji oraz rozliczanie zakładów. Silnik domenowy nie zależy
bezpośrednio od biblioteki Gymnasium. Adapter zgodny z jej interfejsem zostanie
utworzony jako osobna warstwa, co umożliwi późniejsze porównywanie różnych
reprezentacji stanu i funkcji nagrody bez modyfikowania zasad gry.

\section{Założenia projektowe środowiska}
\label{sec:zalozenia-srodowiska}

Podstawową jednostką interakcji ze środowiskiem jest jedno kompletne rozdanie.
Rozdanie rozpoczyna się od utworzenia talii i przekazania graczowi oraz
krupierowi po dwóch kart własnych, a kończy się spasowaniem gracza albo
automatycznym rozstrzygnięciem układów. Z punktu widzenia uczenia przez
wzmacnianie pojedyncze rozdanie stanowi jeden epizod.

Agent podejmuje decyzje wyłącznie w trzech etapach: przed flopem, po odkryciu
flopa oraz po odkryciu wszystkich pięciu kart wspólnych. W zależności od etapu
może czekać, wykonać zakład Play o odpowiednim mnożniku albo spasować.
Po wykonaniu zakładu Play środowisko automatycznie odkrywa pozostałe karty,
ocenia ręce gracza i krupiera oraz przechodzi do stanu końcowego. W ramach
jednego rozdania zakład Play może zostać wykonany tylko raz.

Środowisko modeluje podstawowy wariant Ultimate Texas Hold’em obejmujący
zakłady Ante, Blind oraz Play. Zakłady boczne Trips i Progressive zostały
pominięte. Nie wpływają one na przestrzeń decyzji podstawowej gry, natomiast
zmieniają rozkład wyników finansowych i utrudniają ocenę jakości strategii
decyzyjnej agenta. Uzasadnienie tego ograniczenia przedstawiono szerzej
w dalszej części rozdziału.

Wartości zakładów wyrażono w jednostkach względnych, gdzie jedna jednostka
odpowiada wartości zakładu Ante. Zakłady Ante i Blind mają zatem wartość
jednej jednostki, natomiast wartość zakładu Play zależy od momentu jego
wykonania i może wynosić od jednej do czterech jednostek. Takie rozwiązanie
uniezależnia środowisko od konkretnej waluty i nominalnej wysokości stawki,
a także ułatwia porównywanie wyników pomiędzy eksperymentami.

Silnik zwraca szczegółowe rozliczenie poszczególnych zakładów oraz łączny
wynik netto rozdania. W podstawowym wariancie środowiska przeznaczonym dla
agentów uczących się wynik netto będzie wykorzystywany jako nagroda terminalna.
W krokach pośrednich nagroda wynosi zero. Rozdzielenie mechanizmu rozliczania
gry od funkcji nagrody umożliwia późniejsze badanie alternatywnych wariantów
nagrody bez ingerencji w reguły Ultimate Texas Hold’em.

Istotnym założeniem była również powtarzalność eksperymentów. Standardowa
talia może być tasowana przy użyciu zadanego ziarna generatora liczb
pseudolosowych, dzięki czemu możliwe jest odtworzenie całej sekwencji rozdań.
W testach jednostkowych stosowana jest dodatkowo talia deterministyczna,
pozwalająca jednoznacznie określić kolejność dobieranych kart i zweryfikować
konkretne przypadki przebiegu oraz rozliczenia rozdania.

\section{Architektura rozwiązania}
\label{sec:architektura-srodowiska}

Środowisko Ultimate Texas Hold’em zaprojektowano jako zestaw niezależnych
modułów domenowych. Każdy moduł odpowiada za jeden wyraźnie określony obszar,
taki jak reprezentacja kart, ocena układów, sterowanie przebiegiem rozdania
albo rozliczanie zakładów. Centralnym elementem rozwiązania jest klasa
\texttt{UTHGame}, która koordynuje działanie pozostałych komponentów, ale nie
implementuje bezpośrednio wszystkich reguł gry.

Przyjęty podział ogranicza zależności pomiędzy poszczególnymi częściami
systemu. Model kart i evaluator układów mogą być rozwijane oraz testowane
niezależnie od zasad Ultimate Texas Hold’em. Analogicznie mechanizm
rozliczania zakładów nie odpowiada za odkrywanie kart ani za walidację
dozwolonych akcji. Dzięki temu zmiana jednego elementu, na przykład
reprezentacji obserwacji agenta, nie wymaga modyfikowania sposobu oceny
układów lub obliczania wypłat.

Najniższą warstwę rozwiązania stanowią moduły \texttt{cards} i
\texttt{hands}. Pierwszy z nich definiuje karty oraz talię, natomiast drugi
odpowiada za ocenę i porównywanie układów pokerowych. Moduły te nie zawierają
reguł charakterystycznych dla Ultimate Texas Hold’em i mogą zostać ponownie
wykorzystane w innych odmianach pokera.

Warstwę domenową gry tworzy pakiet \texttt{uth}. Zawiera on modele stanu,
reguły legalności akcji, mechanizm rozdawania kart, sposób rozliczania
zakładów oraz logikę sterującą kolejnymi fazami rozdania. Klasa
\texttt{UTHGame} pełni rolę koordynatora tej warstwy. Na podstawie aktualnego
stanu i wybranej akcji wywołuje odpowiednie operacje, aktualizuje stan
rozdania oraz zwraca rezultat kroku.

Agent nie komunikuje się bezpośrednio z talią, evaluatorem ani pełnym stanem
wewnętrznym gry. Otrzymuje obiekt \texttt{UTHObservation}, a swoją decyzję
przekazuje do metody \texttt{step}. Rezultatem wykonania akcji jest obiekt
\texttt{StepResult}, zawierający kolejną obserwację oraz, w przypadku
zakończenia rozdania, jego wynik i szczegółowe rozliczenie. Takie rozwiązanie
tworzy jednoznaczną granicę pomiędzy logiką środowiska a implementacją agenta.

Tabela~\ref{tab:moduly-srodowiska} przedstawia najważniejsze moduły
rozwiązania i zakres ich odpowiedzialności.

\begin{table}[htbp]
    \centering
    \caption{Podział odpowiedzialności pomiędzy modułami środowiska UTH}
    \label{tab:moduly-srodowiska}
    \begin{tabular}{|p{0.27\textwidth}|p{0.63\textwidth}|}
        \hline
        \textbf{Moduł} & \textbf{Odpowiedzialność} \\
        \hline
        \texttt{cards}
        &
        Reprezentacja rang, kolorów, pojedynczych kart oraz standardowej
        talii składającej się z 52 unikalnych kart.
        \\
        \hline

        \texttt{hands}
        &
        Ocena pięciokartowych układów, wybór najlepszego układu z pięciu,
        sześciu lub siedmiu kart oraz porównywanie wartości rąk.
        \\
        \hline

        \texttt{uth.models}
        &
        Definicja pełnego stanu rozdania, obserwacji agenta oraz wyniku
        pojedynczego kroku środowiska.
        \\
        \hline

        \texttt{uth.rules}
        &
        Określanie zbioru legalnych akcji dla każdej fazy rozdania.
        \\
        \hline

        \texttt{uth.dealing}
        &
        Realizacja kolejności rozdawania kart, odkrywania flopa, turna
        i rivera oraz obsługa kart spalonych.
        \\
        \hline

        \texttt{uth.card\_source}
        &
        Abstrakcja źródła kart oraz deterministyczna talia wykorzystywana
        w testach.
        \\
        \hline

        \texttt{uth.settlement}
        &
        Określanie wyniku rozdania, kwalifikacji krupiera oraz rozliczanie
        zakładów Ante, Blind i Play.
        \\
        \hline

        \texttt{uth.game}
        &
        Koordynacja przebiegu rozdania, walidacja akcji, przechodzenie
        pomiędzy fazami oraz tworzenie wyników kolejnych kroków.
        \\
        \hline

        \texttt{uth.trace}
        &
        Opcjonalne rejestrowanie obserwacji, decyzji agenta i końcowego
        stanu rozdania na potrzeby diagnostyki.
        \\
        \hline
    \end{tabular}
\end{table}

Silnik domenowy nie zależy od konkretnej biblioteki uczenia przez
wzmacnianie. Przyszły adapter Gymnasium będzie korzystał z publicznego
interfejsu \texttt{UTHGame}, przekształcając obiekty domenowe na numeryczne
reprezentacje wymagane przez algorytmy uczące się. Pozwala to zachować jedną
implementację zasad gry dla wszystkich badanych reprezentacji obserwacji,
funkcji nagrody i metod uczenia.